\section{Exact settlement on a fixed region}
\label{sec:settlement}

Fix a load $\ell\in\mathbb R^V$ and define the quadratic
\begin{equation}
 \Psi_\ell(x):=\frac12x^TQx-\ell^Tx,
 \qquad
 e_\ell(x):=\ell-Qx.
 \label{eq:load-quadratic}
\end{equation}
For $U\subseteq V$, let $\iota_U$ denote zero extension from
$\mathbb R^U$ to $\mathbb R^V$.

\begin{definition}[Dirichlet settlement]
\label{def:dirichlet-settlement}
The exact settlement of $x$ on $U$ is
\begin{equation}
 \mathcal E_U^\ell(x)
 :=x+\iota_UQ_{UU}^{-1}e_\ell(x)_U.
 \label{eq:dirichlet-settlement}
\end{equation}
Equivalently, it is the unique minimizer of $\Psi_\ell$ over all $z$ with
$z_{U^c}=x_{U^c}$.
\end{definition}

The definition is always valid because $Q\succeq\alpha I$ makes every
principal block positive definite.

\begin{lemma}[Canonical form and exact energy drop]
\label{lem:canonical-settlement}
For every $x$ and $U$,
\begin{align}
 \mathcal E_U^\ell(x)_{U^c}&=x_{U^c},
 \label{eq:settlement-outside}\\
 \mathcal E_U^\ell(x)_U
 &=Q_{UU}^{-1}\bigl(\ell_U-Q_{UU^c}x_{U^c}\bigr),
 \label{eq:settlement-canonical}\\
 e_\ell(\mathcal E_U^\ell(x))_U&=0,
 \label{eq:settlement-residual-zero}\\
 \Psi_\ell(x)-\Psi_\ell(\mathcal E_U^\ell(x))
 &=\frac12 e_\ell(x)_U^TQ_{UU}^{-1}e_\ell(x)_U.
 \label{eq:settlement-energy-drop}
\end{align}
In particular, the settled point depends on $x$ only through the outside
coordinates $x_{U^c}$.
\end{lemma}

\begin{proof}
Put $z_U=Q_{UU}^{-1}e_\ell(x)_U$.  The definition changes no outside
coordinate, and
\[
 x_U+z_U
 =x_U+Q_{UU}^{-1}
   (\ell_U-Q_{UU}x_U-Q_{UU^c}x_{U^c}),
\]
which proves \cref{eq:settlement-canonical}.  Multiplying that identity by
$Q_{UU}$ proves \cref{eq:settlement-residual-zero}.  Finally, the exact
quadratic expansion along $\iota_Uz_U$ gives
\[
 \Psi_\ell(x+\iota_Uz_U)-\Psi_\ell(x)
 =-e_\ell(x)_U^Tz_U+\frac12z_U^TQ_{UU}z_U,
\]
which is minus the right side of \cref{eq:settlement-energy-drop}.
\end{proof}

\begin{theorem}[Settlement absorption law]
\label{thm:settlement-absorption}
For every fixed load $\ell$:
\begin{enumerate}[label=(\roman*),leftmargin=*]
 \item $\mathcal E_U^\ell\mathcal E_U^\ell=\mathcal E_U^\ell$.
 \item If $U\subseteq T$, then
 \begin{equation}
  \mathcal E_T^\ell\mathcal E_U^\ell
  =\mathcal E_U^\ell\mathcal E_T^\ell
  =\mathcal E_T^\ell.
  \label{eq:nested-absorption}
 \end{equation}
 \item If $y-x$ is supported in $U$, then
 \begin{equation}
  \mathcal E_U^\ell(y)=\mathcal E_U^\ell(x).
  \label{eq:trajectory-erasure}
 \end{equation}
\end{enumerate}
\end{theorem}

\begin{proof}
After one settlement the residual on $U$ is zero, so applying the same map
again changes nothing.  If $U\subseteq T$, settling $U$ changes no coordinate
outside $T$.  The canonical formula \cref{eq:settlement-canonical} says that
settlement on $T$ depends only on those unchanged outside coordinates, which
proves $\mathcal E_T^\ell\mathcal E_U^\ell=\mathcal E_T^\ell$.
Conversely, the $T$-settled point has zero residual throughout $T$, hence
also on $U$, so the later $U$-settlement does nothing.  The last statement is
the same canonical-form argument: $x$ and $y$ agree on $U^c$.
\end{proof}

The theorem is the precise sense in which settlement forgets propagation.
The propagation may be over-relaxed, ranked, asynchronous, or represented as
signed directed packets.  If all its coordinate changes remain inside the
region that is subsequently settled, none of those choices changes the
settled point.

\begin{corollary}[Immediate and block cleanup are one phenomenon]
\label{cor:cleanup-unification}
An over-relaxed coordinate push at $u$ followed immediately by exact
settlement on $\{u\}$ equals one unit-relaxation push.  More generally, every
finite sequence of coordinate updates supported in $U$, followed by exact
settlement on $U$, equals direct settlement of the pre-sequence iterate.
\end{corollary}

\begin{proof}
Apply \cref{eq:trajectory-erasure} with $U=\{u\}$ or with the full update
support $U$.
\end{proof}
